\section{Reproducibilidad y Entornos de Ejecución}
\label{sec:reproducibilidad}

La reproducibilidad científica fue un principio rector en el diseño de ASDAT-253. Para garantizar que cualquier investigador o evaluador académico pueda replicar la totalidad de los experimentos, transformaciones y reportes, el proyecto implementa un protocolo determinista basado en aislamiento de dependencias, semillas pseudoaleatorias fijas y trazabilidad transaccional.

\subsection{Gestión Dual de Entornos: R y Python}

El sistema desacopla sus dos núcleos computacionales mediante gestores de paquetes reproducibles:
\begin{itemize}
  \item \textbf{Ecosistema R (Investigación y Modelado Clásico):} Las librerías estadísticas (\texttt{quanteda}, \texttt{topicmodels}, \texttt{testthat}) se congelan mediante \texttt{renv} con el archivo de bloqueo \texttt{renv.lock} (248 KB con hashes de integridad).
  \item \textbf{Ecosistema Python (Producción Web y Paridad):} El microframework Flask, los adaptadores \texttt{psycopg2-binary}, \texttt{pandas} y la suite \texttt{pytest} se gestionan mediante \texttt{requirements.txt} sobre Python 3.12.
\end{itemize}

\subsection{Pipeline de Ejecución Secuencial}

La pipeline de investigación en R se ejecuta de forma determinista mediante scripts numerados:
\begin{verbatim}
Rscript scripts/00_run_all.R         # Orquestador V1 y V2 completo
Rscript scripts/10_test.R            # Suite de pruebas unitarias testthat
\end{verbatim}

De forma paralela, la suite de paridad matemática y los servicios web en Python se ejecutan mediante:
\begin{verbatim}
python -m pytest tests/test_emotica_parity.py -v  # Paridad R-Python
python -m pytest tests/test_ingest_endpoint.py -v# Ingesta multipart/ZIP
python run.py                                    # Servidor Flask (puerto 5000)
\end{verbatim}

\subsection{Persistencia y Trazabilidad en Neon PostgreSQL}

El aprovisionamiento de la base de datos distribuida en la nube se ejecuta de forma automatizada mediante scripts idempotentes en \texttt{database/}:
\begin{enumerate}
  \item \texttt{neon\_complete\_schema.sql}: DDL maestro con las 16 tablas, índices y triggers CDC.
  \item \texttt{migrate\_all\_to\_neon.py}: Migración masiva de 23.371 mensajes, matrices TF-IDF y tópicos.
  \item \texttt{sync\_catalogos\_umsa.py}: Sincronización de docentes, materias y oferta semestral.
  \item \texttt{detectar\_menciones\_umsa.py}: Escaneo y resolución de entidades con expresiones regulares.
\end{enumerate}

\subsection{Generación Automatizada del Informe}

Ninguna cifra en este informe ha sido digitada manualmente. Todas las macros estadísticas (\texttt{\textbackslash nMensajes}, \texttt{\textbackslash pctCobertura}, \texttt{\textbackslash piNeutro}, etc.) y las tablas formales se regeneran automáticamente desde las salidas analíticas en \texttt{reports/generated/values.tex} y \texttt{reports/tables/}.
